\section{Conclusion}
\label{sec:conclusion}

The central finding of this paper is that the gain attributed to LLM-guided memory fusion in fact-centric long-horizon conversational QA comes almost entirely from deduplicating and packing verbatim evidence into a fixed retrieval budget --- not from the LLM's linguistic rewriting ability.
DFM-Fusion, a fully deterministic four-stage operator requiring zero maintenance-time LLM calls, recovers this gain across a heterogeneous family of six open-weight models.
DFM-Route, its query-aware retrieval extension, further demonstrates that the remaining temporal gap is a retrieval-routing problem rather than a rewriting problem: by routing temporal queries to anchor-based retrieval with date-unit recovery, temporal F1 improves $+57\%$ (from 0.190 to 0.299) without any LLM call at maintenance or retrieval time.
Together, DFM-Fusion and DFM-Route demonstrate that the full memory pipeline --- from consolidation to evidence delivery --- can operate without any API dependency in the maintenance loop.

At the same time, our results do not support a blanket equivalence claim.
DFM-Fusion's core advantage is specific to fact-centric QA, where the evaluation rewards verbatim evidence retrieval; open-domain preference synthesis may still benefit from paraphrastic rewriting.
An adversarial trade-off accompanies the expanded retrieval in DFM-Route: routing more queries to union or expansive strategies introduces noise that adversarial questions penalize, a tension that warrants further study.
Future work should extend this investigation to benchmarks requiring abstraction and explicit conflict reconciliation, and explore hybrid operators that selectively invoke LLM rewriting only when coverage verification signals a genuine synthesis gap.
Limitations of the current study --- including single-benchmark scope, surface-biased preservation heuristics, and the absence of a held-out validation split --- are discussed in Section~\ref{sec:discussion}.
